\chapter{Fundamentos del Electromagnetismo}
\label{cap:tema1}

\cuadrofuentes{
    \textbf{Asignatura:} Fundamentos Físicos y Tecnológicos (1.º Grado en Ingeniería Informática, UGR).\\[2pt]
    \textbf{Profesor:} Pedro Cartujo (Departamento de Electrónica y Tecnología de Computadores).\\[2pt]
    \textbf{Bibliografía de Referencia:} Sears, F. W., Zemansky, M. W., Young, H. D., Freedman, R. A., \textit{Física Universitaria con Física Moderna} (Vol. 2), Addison-Wesley.
}

% ==============================================================================
% 1.1. ELECTROSTÁTICA
% ==============================================================================
\section{Electrostática: Carga Eléctrica, Fuerzas y Campo Eléctrico}

El electromagnetismo unifica los fenómenos eléctricos y magnéticos bajo un marco físico común. La electrostática se ocupa de las interacciones producidas por cargas eléctricas que se encuentran en reposo relativo (campo electrostático independiente del tiempo).

\subsection{Concepto y Propiedades de la Carga Eléctrica}
La carga eléctrica es una propiedad intrínseca y fundamental de la materia que determina su interacción electromagnética. Existen dos tipos de carga: positiva (asociada a los protones) y negativa (asociada a los electrones).

\begin{definicion}[Propiedades Universales de la Carga Eléctrica]
\begin{itemize}[leftmargin=*]
    \item \textbf{Ley de Conservación de la Carga:} En cualquier proceso físico que tenga lugar dentro de un sistema eléctricamente aislado, la carga eléctrica total neta permanece rigurosamente constante ($\sum q = \text{constante}$). La carga no se crea ni se destruye, únicamente se transfiere o redistribuye.
    \item \textbf{Cuantización de la Carga:} La carga neta de cualquier cuerpo macroscópico o partícula subatómica observable es siempre un múltiplo entero de la carga elemental fundamental $e$:
    \begin{equation}
        q = \pm n \cdot e, \quad \text{con } n \in \mathbb{N}, \quad e \approx 1{,}602 \times 10^{-19}\ \si{\coulomb}
    \end{equation}
    \item \textbf{Invarianza Relativista:} La magnitud de la carga eléctrica de una partícula es independiente de su velocidad o del sistema de referencia inercial desde el cual se observe.
\end{itemize}
\end{definicion}

En el Sistema Internacional (SI), la unidad de carga eléctrica es el \textbf{Culombio} (\si{\coulomb}), definido como la cantidad de carga que atraviesa la sección transversal de un conductor por el que circula una corriente constante de un Amperio durante un segundo ($1\ \si{\coulomb} = 1\ \si{\ampere\second}$).

\subsection{Ley de Coulomb}
En 1785, Charles-Augustin de Coulomb determinó experimentalmente la fuerza electrostática entre dos cargas puntuales en reposo en el vacío.

\begin{ley}[Ley de Coulomb]
La fuerza electrostática $\vec{F}_{12}$ que una carga puntual $q_1$ ejerce sobre otra carga puntual $q_2$ separadas por una distancia $r$ es directamente proporcional al producto de sus cargas e inversamente proporcional al cuadrado de la distancia que las separa:
\begin{equation}
    \vec{F}_{12} = \frac{1}{4\pi\varepsilon_0} \frac{q_1 q_2}{r^2} \hat{r}_{12} = K_e \frac{q_1 q_2}{r^2} \hat{r}_{12}
\end{equation}
donde:
\begin{itemize}[leftmargin=*]
    \item $\hat{r}_{12} = \frac{\vec{r}_2 - \vec{r}_1}{|\vec{r}_2 - \vec{r}_1|}$ es el vector unitario dirigido desde la carga fuente $q_1$ hacia la carga testigo $q_2$.
    \item $\varepsilon_0 \approx 8{,}854 \times 10^{-12}\ \si{\coulomb\squared\per\newton\per\meter\squared}$ es la permitividad dieléctrica del vacío.
    \item $K_e = \frac{1}{4\pi\varepsilon_0} \approx 8{,}988 \times 10^9\ \si{\newton\meter\squared\per\coulomb\squared} \approx 9 \times 10^9\ \si{\newton\meter\squared\per\coulomb\squared}$.
\end{itemize}
\end{ley}

\begin{center}
\begin{tikzpicture}[scale=1.1]
    % Carga q1
    \draw[fill=fftPrimary!20, draw=fftPrimary, thick] (-2.5,0) circle (0.35) node[fftPrimary] {\textbf{$+q_1$}};
    % Carga q2
    \draw[fill=fftPrimary!20, draw=fftPrimary, thick] (2.5,0) circle (0.35) node[fftPrimary] {\textbf{$+q_2$}};
    % Linea de separación
    \draw[dashed, gray] (-2.15,0) -- (2.15,0) node[midway, below=2pt, text=black] {$r$};
    % Fuerzas repulsivas
    \draw[->, >=stealth, very thick, fftAccent] (-2.85,0) -- (-4.2,0) node[above] {$\vec{F}_{21}$};
    \draw[->, >=stealth, very thick, fftAccent] (2.85,0) -- (4.2,0) node[above] {$\vec{F}_{12}$};
    \node[below=14pt] at (0,0) {\small\textsf{Fuerzas electrostáticas repulsivas entre dos cargas del mismo signo ($\vec{F}_{21} = -\vec{F}_{12}$)}};
\end{tikzpicture}
\end{center}

\subsection{Campo Eléctrico y Principio de Superposición}
Para evitar el concepto de acción a distancia instantánea, se define el campo eléctrico como una perturbación vectorial del espacio circundante provocada por la presencia de cargas eléctricas.

\begin{definicion}[Vector Campo Eléctrico $\vec{E}$ y Líneas de Campo]
El campo eléctrico $\vec{E}$ en un punto $\vec{r}$ del espacio es la fuerza electrostática por unidad de carga positiva testigo $q_0$ situada en dicho punto:
\begin{equation}
    \vec{E}(\vec{r}) = \lim_{q_0 \to 0} \frac{\vec{F}}{q_0} \quad \left[\si{\newton\per\coulomb} \equiv \si{\volt\per\meter}\right]
\end{equation}
Para una carga puntual aislada:
\begin{equation}
    \vec{E}(\vec{r}) = \frac{1}{4\pi\varepsilon_0} \frac{q}{r^2} \hat{r}
\end{equation}
Las líneas de campo eléctrico parten de las cargas positivas (fuentes) y convergen hacia las cargas negativas (sumideros).
\end{definicion}

\begin{center}
\begin{tikzpicture}[scale=0.9]
    % Carga positiva
    \draw[fill=fftGreen!25, draw=fftGreen, thick] (-3.5,0) circle (0.3) node[fftGreen] {\textbf{$+$}};
    \foreach \angle in {0, 45, 90, 135, 180, 225, 270, 315} {
        \draw[->, >=stealth, thick, fftGreen] ({-3.5 + 0.35*cos(\angle)}, {0.35*sin(\angle)}) -- ({-3.5 + 1.2*cos(\angle)}, {1.2*sin(\angle)});
    }
    \node[below=12pt] at (-3.5,-1) {\footnotesize\textsf{Carga positiva (fuente)}};

    % Carga negativa
    \draw[fill=fftRed!25, draw=fftRed, thick] (3.5,0) circle (0.3) node[fftRed] {\textbf{$-$}};
    \foreach \angle in {0, 45, 90, 135, 180, 225, 270, 315} {
        \draw[<-, >=stealth, thick, fftRed] ({3.5 + 0.35*cos(\angle)}, {0.35*sin(\angle)}) -- ({3.5 + 1.2*cos(\angle)}, {1.2*sin(\angle)});
    }
    \node[below=12pt] at (3.5,-1) {\footnotesize\textsf{Carga negativa (sumidero)}};
\end{tikzpicture}
\end{center}

\begin{ley}[Principio de Superposición Lineal]
El campo eléctrico neto generado por una distribución de $N$ cargas puntuales $\{q_1, q_2, \dots, q_N\}$ es la suma vectorial de los campos individuales generados por cada una de ellas:
\begin{equation}
    \vec{E}(\vec{r}) = \sum_{i=1}^N \vec{E}_i(\vec{r}) = \frac{1}{4\pi\varepsilon_0} \sum_{i=1}^N \frac{q_i}{|\vec{r} - \vec{r}_i|^2} \frac{\vec{r} - \vec{r}_i}{|\vec{r} - \vec{r}_i|}
\end{equation}
Para distribuciones continuas de carga, se integra sobre la densidad correspondiente:
\begin{itemize}[leftmargin=*]
    \item Densidad lineal de carga: $\lambda = \frac{dq}{d\ell}\quad [\si{\coulomb\per\meter}] \implies dq = \lambda\, d\ell$.
    \item Densidad superficial de carga: $\sigma = \frac{dq}{dA}\quad [\si{\coulomb\per\meter\squared}] \implies dq = \sigma\, dA$.
    \item Densidad volumétrica de carga: $\rho = \frac{dq}{dV}\quad [\si{\coulomb\per\meter\cubed}] \implies dq = \rho\, dV$.
\end{itemize}
\end{ley}

% ==============================================================================
% 1.2. LEY DE GAUSS
% ==============================================================================
\section{Flujo Eléctrico y Ley de Gauss}

\subsection{Concepto de Flujo Eléctrico y Ley de Gauss}
\begin{ley}[Ley de Gauss]
El flujo del campo eléctrico a través de cualquier superficie gaussiana cerrada $S_{cerrada}$ es directamente proporcional a la carga eléctrica neta total $Q_{enc}$ encerrada en el interior de dicha superficie:
\begin{equation}
    \Phi_E = \oint_S \vec{E} \cdot d\vec{A} = \frac{Q_{enc}}{\varepsilon_0}
\end{equation}
\end{ley}

\begin{center}
\begin{tikzpicture}[scale=1.0]
    % Hilo cargado con lambda
    \draw[very thick, fftAccent] (-3,0) -- (3,0) node[right] {Hilo indefinido ($+\lambda$)};
    \foreach \x in {-2.5,-1.8,...,2.5} {
        \node[above=2pt, fftAccent, font=\tiny] at (\x,0) {$+$};
    }
    % Superficie gaussiana cilíndrica
    \draw[thick, fftPrimary, dashed] (-1.8,-1) rectangle (1.8,1);
    \draw[thick, fftPrimary] (-1.8,0) ellipse (0.25 and 1);
    \draw[thick, fftPrimary] (1.8,0) ellipse (0.25 and 1);
    % Vectores campo y área
    \draw[->, >=stealth, very thick, fftGreen] (0,1) -- (0,2.1) node[above] {$\vec{E}$};
    \draw[->, >=stealth, very thick, fftDarkBlue] (0,1) -- (0,1.6) node[right] {$d\vec{A}$};
    \draw[<->, gray] (1.9,0) -- (1.9,1) node[midway, right] {$r$};
    \draw[<->, gray] (-1.8,-1.2) -- (1.8,-1.2) node[midway, below] {$L$};
    \node[below=16pt] at (0,-1.2) {\small\textsf{Superficie gaussiana cilíndrica coaxial para calcular el campo de un hilo indefinido}};
\end{tikzpicture}
\end{center}

\subsection{Conductores en Equilibrio Electrostático y Condensadores}
\begin{definicion}[Condensador de Placas Paralelas]
Dispositivo formado por dos láminas conductoras paralelas de área $A$ separadas por una distancia $d$ que almacenan cargas opuestas $+Q$ y $-Q$.
\begin{equation}
    C = \frac{Q}{\Delta V} = \varepsilon_r \varepsilon_0 \frac{A}{d} \quad [\si{\farad}]
\end{equation}
\end{definicion}

\begin{center}
\begin{tikzpicture}[scale=1.1]
    % Placa positiva
    \draw[fill=fftPrimary!20, thick] (-2,1) rectangle (2,1.3);
    \node at (0,1.15) {\small\textbf{Placa Superior ($+Q$, densidad $+\sigma$)}};
    % Placa negativa
    \draw[fill=fftDarkBlue!20, thick] (-2,-1.3) rectangle (2,-1);
    \node at (0,-1.15) {\small\textbf{Placa Inferior ($-Q$, densidad $-\sigma$)}};
    % Líneas de campo uniforme
    \foreach \x in {-1.5,-1.0,-0.5,0,0.5,1.0,1.5} {
        \draw[->, >=stealth, thick, fftAccent] (\x,1) -- (\x,-1);
    }
    \node[right, fftAccent] at (1.6,0) {$\vec{E} = \frac{\sigma}{\varepsilon}$};
    % Cota de separación
    \draw[<->, gray, thick] (-2.3,-1) -- (-2.3,1) node[midway, left] {$d$};
    % Conexión exterior
    \draw[thick] (0,1.3) -- (0,1.8) node[above] {$V_+$};
    \draw[thick] (0,-1.3) -- (0,-1.8) node[below] {$V_-$};
\end{tikzpicture}
\end{center}

\subsection{Asociación de Condensadores}
\begin{center}
\begin{tikzpicture}[scale=0.9]
    % En paralelo
    \node[font=\bfseries] at (-3,1.8) {Asociación en Paralelo ($C_{eq} = C_1 + C_2$)};
    \draw (-4.5,0) -- (-3.5,0);
    \draw (-3.5,0) -- (-3.5,0.8) to[C, l=$C_1$] (-1.5,0.8) -- (-1.5,0);
    \draw (-3.5,0) -- (-3.5,-0.8) to[C, l=$C_2$] (-1.5,-0.8) -- (-1.5,0);
    \draw (-1.5,0) -- (-0.5,0);

    % En serie
    \node[font=\bfseries] at (3.5,1.8) {Asociación en Serie ($\frac{1}{C_{eq}} = \frac{1}{C_1} + \frac{1}{C_2}$)};
    \draw (1.5,0) -- (2,0) to[C, l=$C_1$] (3.5,0) to[C, l=$C_2$] (5,0) -- (5.5,0);
\end{tikzpicture}
\end{center}

% ==============================================================================
% 1.3. CAMPO MAGNÉTICO Y LEYES FUNDAMENTALES
% ==============================================================================
\section{Campo Magnético, Ley de Lorentz y Ley de Biot-Savart}

\begin{ley}[Fuerza de Lorentz y Ley de Biot-Savart]
Una partícula con carga $q$ que se mueve con velocidad $\vec{v}$ en un campo $\vec{B}$ experimenta la fuerza:
\begin{equation}
    \vec{F}_B = q (\vec{v} \times \vec{B}) \implies |\vec{F}_B| = |q|\, v\, B \sin(\theta)
\end{equation}
El campo magnético generado por un conductor rectilíneo muy largo por el que circula una corriente $I$ viene dado por la Ley de Biot-Savart / Ampère:
\begin{equation}
    B(r) = \frac{\mu_0 I}{2\pi r}
\end{equation}
\end{ley}

\begin{center}
\begin{tikzpicture}[scale=1.0]
    % Hilo con corriente
    \draw[very thick, fftDarkBlue] (0,-1.8) -- (0,1.8) node[above] {$I$ (corriente)};
    \draw[->, >=stealth, very thick, fftDarkBlue] (0,-0.5) -- (0,0.6);
    % Líneas de campo circulares
    \draw[thick, fftAccent, ->, >=stealth] (0.8,0) arc (0:170:0.8 and 0.35);
    \draw[thick, fftAccent, dashed] (0.8,0) arc (0:-190:0.8 and 0.35);
    \draw[thick, fftAccent, ->, >=stealth] (1.5,0) arc (0:170:1.5 and 0.6);
    \draw[thick, fftAccent, dashed] (1.5,0) arc (0:-190:1.5 and 0.6);
    \node[right, fftAccent] at (1.5,0.4) {$\vec{B}$ (Líneas concéntricas)};
    \node[below=14pt] at (0,-1.8) {\small\textsf{Líneas de campo magnético según la Regla de la Mano Derecha}};
\end{tikzpicture}
\end{center}

\begin{ley}[Fuerza entre Dos Conductores Paralelos]
Para dos conductores rectilíneos paralelos separados una distancia $d$:
\begin{equation}
    \frac{F}{\ell} = \frac{\mu_0 I_1 I_2}{2\pi d}
\end{equation}
\begin{itemize}[leftmargin=*]
    \item Corrientes en el \textbf{mismo sentido}: fuerzas de \textbf{atracción mutua}.
    \item Corrientes en \textbf{sentidos opuestos}: fuerzas de \textbf{repulsión mutua}.
\end{itemize}
\end{ley}

\begin{center}
\begin{tikzpicture}[scale=1.0]
    % Caso 1: Corrientes mismo sentido
    \draw[very thick, fftPrimary] (-3.5,-1.2) -- (-3.5,1.2) node[above] {$I_1$};
    \draw[->, >=stealth, very thick, fftPrimary] (-3.5,-0.4) -- (-3.5,0.4);
    \draw[very thick, fftPrimary] (-1.5,-1.2) -- (-1.5,1.2) node[above] {$I_2$};
    \draw[->, >=stealth, very thick, fftPrimary] (-1.5,-0.4) -- (-1.5,0.4);
    \draw[->, >=stealth, very thick, fftGreen] (-3.5,0) -- (-2.4,0) node[midway, above] {$\vec{F}_{21}$};
    \draw[->, >=stealth, very thick, fftGreen] (-1.5,0) -- (-2.6,0) node[midway, above] {$\vec{F}_{12}$};
    \node[below=8pt] at (-2.5,-1.2) {\small\textbf{Atracción} (mismo sentido)};

    % Caso 2: Corrientes sentido opuesto
    \draw[very thick, fftDarkBlue] (1.5,-1.2) -- (1.5,1.2) node[above] {$I_1$};
    \draw[->, >=stealth, very thick, fftDarkBlue] (1.5,-0.4) -- (1.5,0.4);
    \draw[very thick, fftDarkBlue] (3.5,-1.2) -- (3.5,1.2) node[above] {$I_2$};
    \draw[<-, >=stealth, very thick, fftDarkBlue] (3.5,-0.4) -- (3.5,0.4);
    \draw[->, >=stealth, very thick, fftRed] (1.5,0) -- (0.6,0) node[midway, above] {$\vec{F}$};
    \draw[->, >=stealth, very thick, fftRed] (3.5,0) -- (4.4,0) node[midway, above] {$\vec{F}$};
    \node[below=8pt] at (2.5,-1.2) {\small\textbf{Repulsión} (sentidos opuestos)};
\end{tikzpicture}
\end{center}

\subsection{Inducción Electromagnética y Autoinducción}
\begin{ley}[Ley de Faraday-Lenz y Autoinducción en Solenoide]
La fuerza electromotriz inducida en un circuito cerrado de $N$ espiras ante una variación de flujo magnético es:
\begin{equation}
    \mathcal{E} = -N \frac{d\Phi_B}{dt}
\end{equation}
Para un solenoide de longitud $\ell$, área $A$ y $N$ espiras, el coeficiente de autoinducción es:
\begin{equation}
    L = \frac{N \Phi_B}{I} = \mu_0 \frac{N^2 A}{\ell} \quad [\si{\henry}]
\end{equation}
La tensión autoinducida al variar la corriente es $\mathcal{E}_L = -L \frac{di}{dt}$, almacenando la energía magnética $U_B = \frac{1}{2} L I^2$.
\end{ley}

% ==============================================================================
% 1.4. EJERCICIOS RESUELTOS DE EXAMEN
% ==============================================================================
\section{Ejercicios Resueltos de Exámenes de la UGR}

\begin{ejercicio}[Cálculo de Campo y Potencial de Hilo Finito e Infinito --- Examen Febrero 2017]
Se tiene un hilo rectilíneo de longitud infinita cargado uniformemente con densidad lineal constante $\lambda = 4 \times 10^{-9}\ \si{\coulomb\per\meter}$.
\begin{enumerate}[label=\alph*)]
    \item Calcular el campo eléctrico a una distancia $r_1 = 0{,}05\ \si{\meter}$ del hilo.
    \item Calcular la diferencia de potencial entre los puntos $r_1 = 0{,}05\ \si{\meter}$ y $r_2 = 0{,}20\ \si{\meter}$.
\end{enumerate}
\tcbline
\textbf{Esquema Geométrico del Problema:}\\[4pt]
\begin{center}
\begin{tikzpicture}[scale=1.1]
    % Hilo
    \draw[very thick, fftAccent] (-3,0) -- (3,0) node[right] {Hilo con $\lambda$};
    % Puntos r1 y r2
    \draw[dashed, gray] (0,0) -- (0,2.5);
    \filldraw[fftPrimary] (0,0.8) circle (2pt) node[right=4pt] {$P_1\ (r_1 = 0{,}05\ \si{\meter})$};
    \draw[->, >=stealth, thick, fftGreen] (0,0.8) -- (0,1.7) node[right] {$\vec{E}(r_1)$};
    \filldraw[fftPrimary] (0,2.2) circle (2pt) node[right=4pt] {$P_2\ (r_2 = 0{,}20\ \si{\meter})$};
    \draw[<->, gray] (-0.3,0) -- (-0.3,0.8) node[midway, left] {$r_1$};
    \draw[<->, gray] (-0.7,0) -- (-0.7,2.2) node[midway, left] {$r_2$};
\end{tikzpicture}
\end{center}
\textbf{Solución Paso a Paso:}\\[4pt]
a) Por simetría cilíndrica, aplicando la Ley de Gauss:
\begin{equation}
    E(r) = \frac{\lambda}{2\pi\varepsilon_0 r} = \frac{2 K_e \lambda}{r}
\end{equation}
\begin{equation}
    E(0{,}05) = \frac{2 \times (9 \times 10^9) \times (4 \times 10^{-9})}{0{,}05} = \frac{72}{0{,}05} = 1440\ \si{\volt\per\meter} \quad \implies \vec{E} = 1440\, \hat{u}_r\ \si{\volt\per\meter}
\end{equation}
b) Diferencia de potencial integrando el campo:
\begin{equation}
    V(r_1) - V(r_2) = \frac{\lambda}{2\pi\varepsilon_0} \ln\left(\frac{r_2}{r_1}\right) = 72 \cdot \ln\left(\frac{0{,}20}{0{,}05}\right) = 72 \cdot \ln(4) \approx 99{,}81\ \si{\volt}
\end{equation}
\end{ejercicio}

\begin{ejercicio}[Campo Magnético de Tres Hilos Paralelos --- Examen Enero 2021]
Tres hilos conductores paralelos indefinidos pasan por $x_1 = -d$, $x_2 = 0$ y $x_3 = +d$ ($d = 0{,}10\ \si{\meter}$). Por 1 y 3 circulan $I_1 = 5\ \si{\ampere}$ e $I_3 = 5\ \si{\ampere}$ saliendo ($+\hat{k}$), y por 2 circula $I_2 = 10\ \si{\ampere}$ entrando ($-\hat{k}$).
Determinar la fuerza magnética por unidad de longitud sobre el hilo 2 y sobre el hilo 3.
\tcbline
\textbf{Diagrama de Fuerzas en Corte Transversal:}\\[4pt]
\begin{center}
\begin{tikzpicture}[scale=1.2]
    % Eje horizontal
    \draw[->, gray] (-2.2,0) -- (2.5,0) node[right] {$x$};
    % Hilo 1 (x = -d)
    \draw[thick, fftPrimary] (-1.5,0) circle (0.25);
    \filldraw[fftPrimary] (-1.5,0) circle (2pt);
    \node[below=8pt] at (-1.5,0) {$H_1\ (\odot\ I_1)$};
    % Hilo 2 (x = 0)
    \draw[thick, fftDarkBlue] (0,0) circle (0.25);
    \draw[thick, fftDarkBlue] (-0.17,-0.17) -- (0.17,0.17);
    \draw[thick, fftDarkBlue] (-0.17,0.17) -- (0.17,-0.17);
    \node[below=8pt] at (0,0) {$H_2\ (\otimes\ I_2)$};
    % Hilo 3 (x = +d)
    \draw[thick, fftPrimary] (1.5,0) circle (0.25);
    \filldraw[fftPrimary] (1.5,0) circle (2pt);
    \node[below=8pt] at (1.5,0) {$H_3\ (\odot\ I_3)$};
    % Fuerzas en Hilo 2
    \draw[->, >=stealth, very thick, fftAccent] (0,0.4) -- (0.8,0.4) node[above] {$\vec{F}_{1 \to 2}$};
    \draw[->, >=stealth, very thick, fftGreen] (0,0.4) -- (-0.8,0.4) node[above] {$\vec{F}_{3 \to 2}$};
    % Cotas
    \draw[<->, gray] (-1.5,-0.7) -- (0,-0.7) node[midway, below] {$d$};
    \draw[<->, gray] (0,-0.7) -- (1.5,-0.7) node[midway, below] {$d$};
\end{tikzpicture}
\end{center}
\textbf{Solución Paso a Paso:}\\[4pt]
1) En el hilo central 2: el hilo 1 (corriente contraria) lo repele hacia la derecha ($+\hat{i}$) con $1 \times 10^{-4}\ \si{\newton\per\meter}$; el hilo 3 lo repele hacia la izquierda ($-\hat{i}$) con idéntica magnitud. Por simetría:
\begin{equation}
    \vec{F}_2^{neta} / \ell = \vec{0}
\end{equation}
2) En el hilo 3: el hilo 2 lo repele hacia la derecha ($+1 \times 10^{-4}\ \si{\newton\per\meter}\hat{i}$) y el hilo 1 (corriente en mismo sentido a distancia $2d$) lo atrae hacia la izquierda ($-0{,}25 \times 10^{-4}\ \si{\newton\per\meter}\hat{i}$):
\begin{equation}
    \frac{\vec{F}_3^{neta}}{\ell} = (1{,}00 - 0{,}25) \times 10^{-4}\ \hat{i} = +0{,}75 \times 10^{-4}\ \si{\newton\per\meter}\ \hat{i}
\end{equation}
\end{ejercicio}
